\chapter*{Ek: Cevap Anahtarı}
\addcontentsline{toc}{chapter}{Ek: Cevap Anahtarı}
\markboth{Ek: Cevap Anahtarı}{Ek: Cevap Anahtarı}

\section*{Bölüm 1}
\begin{description}
\item[B1.1 B] Parametre evrene ait bilinmeyen sayısal özelliktir.
\item[B1.2 B] Her satırın temsil ettiği öğrenci analiz birimidir.
\item[B1.3 C] Kimlik numarası sayısal görünse de nominal etikettir.
\item[B1.4 C] Rastgele atama başlangıç özelliklerini dengeleme olasılığını artırır.
\item[B1.5 B] İşlemsel tanım kuramsal yapıyı ölçülebilir hâle getirir.
\item[B1.6 B] Gönüllüler hedef evrenden sistematik biçimde farklı olabilir.
\end{description}

\section*{Bölüm 2}
\begin{description}
\item[B2.1 B] Olasılıklı örneklemede seçilme olasılıkları bilinir.
\item[B2.2 A] Tabakalama önemli alt grupların temsilini güvenceye alır.
\item[B2.3 A] Küme içi benzerlik bağımsız bilgi miktarını azaltır.
\item[B2.4 B] Yanıtlamama sonuçla ilişkiliyse tahmin yanlılaşabilir.
\item[B2.5 A] Ağırlık çoğunlukla seçilme olasılığının tersidir.
\item[B2.6 A] Çerçevedeki eksiklik kapsama hatasıdır.
\end{description}

\section*{Bölüm 3}
\begin{description}
\item[B3.1 B] Medyan uç değerlere karşı ortalamadan daha sağlamdır.
\item[B3.2 A] Varyans standart sapmanın karesidir.
\item[B3.3 A] $IQR=Q_3-Q_1=8$'dir.
\item[B3.4 B] Sabit eklemek yayılmayı değiştirmez.
\item[B3.5 C] Ortalama uç değerlerin büyüklüğünden doğrudan etkilenir.
\item[B3.6 B] Sınıf orta noktası kendi sınıfındaki gözlemleri yaklaşık temsil eder.
\end{description}

\section*{Bölüm 4}
\begin{description}
\item[B4.1 A] $z=1{,}5$, ortalamanın 1,5 standart sapma üzerini gösterir.
\item[B4.2 A] Standart normal dağılımın ortalaması 0, standart sapması 1'dir.
\item[B4.3 B] Merkezî yüzde 95 yaklaşık $\pm1{,}96$ arasındadır.
\item[B4.4 C] Ortalama standart hatası $s/\sqrt n$'dir.
\item[B4.5 B] Örneklem dört katına çıktığında standart hata yarıya iner.
\item[B4.6 B] Teorem örneklem ortalamasının dağılımını açıklar.
\end{description}

\section*{Bölüm 5}
\begin{description}
\item[B5.1 B] Örneklem oranı $x/n$, evren oranını tahmin eder.
\item[B5.2 B] Yansızlık tekrarlı örneklemedeki beklenen değere ilişkindir.
\item[B5.3 A] $MSE=\operatorname{Var}+\operatorname{Bias}^2$'dir.
\item[B5.4 A] Bu özellik tutarlılık olarak adlandırılır.
\item[B5.5 A] Sağlam yöntemler küçük sapma ve uç değerlerden az etkilenir.
\item[B5.6 A] Tabaka payları ağırlıklı tahminde kullanılmalıdır.
\end{description}

\section*{Bölüm 6}
\begin{description}
\item[B6.1 B] Yüzde 95, yöntemin uzun dönem kapsama oranıdır.
\item[B6.2 B] Daha yüksek güven daha büyük kritik değer gerektirir.
\item[B6.3 A] Standart hata azalınca hata payı da azalır.
\item[B6.4 A] Bilinmeyen $\sigma$ için küçük örneklemde $t$ dağılımı kullanılır.
\item[B6.5 B] Sıfır farkı aralıkla uyumlu değerlerden biridir.
\item[B6.6 C] Güven aralığı sistematik seçim yanlılığını kapsamaz.
\end{description}

\section*{Bölüm 7}
\begin{description}
\item[B7.1 B] $p$, $H_0$ altında gözlenen kadar uç sonuç olasılığıdır.
\item[B7.2 A] Doğru sıfır hipotezini reddetmek Tip I hatadır.
\item[B7.3 B] Güç $1-\beta$'dır.
\item[B7.4 B] Yön veri görülmeden önce gerekçelendirilmelidir.
\item[B7.5 A] Büyük örneklem standart hatayı küçültebilir.
\item[B7.6 A] Test sayısı arttıkça yanlış pozitif riski büyür.
\end{description}

\section*{Bölüm 8}
\begin{description}
\item[B8.1 B] Aynı kişilerin ölçümleri eşleştirilmiş yapıdadır.
\item[B8.2 A] Tek örneklem testi ortalamayı sabit değerle karşılaştırır.
\item[B8.3 C] Welch testi eşit evren varyansları gerektirmez.
\item[B8.4 B] Analiz her eşleşmenin fark puanı üzerinden yürür.
\item[B8.5 A] Payda tahminin standart hatasıdır.
\item[B8.6 A] Cohen'in $d$ değeri standartlaştırılmış ortalama farkıdır.
\end{description}

\section*{Bölüm 9}
\begin{description}
\item[B9.1 A] Genel sıfır hipotezi bütün evren ortalamalarının eşitliğidir.
\item[B9.2 A] $F$, gruplar arası ve grup içi ortalama karelerin oranıdır.
\item[B9.3 B] Farkın yeri ayrıca karşılaştırılmalıdır.
\item[B9.4 A] Tukey HSD benzer varyanslı bütün ikilileri denetler.
\item[B9.5 A] Games--Howell eşit olmayan varyanslara uygundur.
\item[B9.6 A] Eta-kare gruplarla ilişkili değişkenlik oranını verir.
\end{description}

\section*{Bölüm 10}
\begin{description}
\item[B10.1 A] Pearson katsayısı doğrusal birlikteliği özetler.
\item[B10.2 B] Sıfır korelasyon eğrisel ilişkiyi dışlamaz.
\item[B10.3 A] Spearman ham değerlerin sıralarını kullanır.
\item[B10.4 A] $(-0{,}70)^2=0{,}49$'dur.
\item[B10.5 A] Ters yön ve ortak nedenler nedensel yorumu sınırlar.
\item[B10.6 A] Saçılım grafiği yön, biçim ve etkili gözlemleri gösterir.
\end{description}

\section*{Bölüm 11}
\begin{description}
\item[B11.1 A] Eğim X'teki bir birime karşı beklenen Y değişimidir.
\item[B11.2 A] Artık $Y-\widehat Y$ olarak hesaplanır.
\item[B11.3 A] En küçük kareler artık kareleri toplamını küçültür.
\item[B11.4 A] Tek açıklayıcılı kesmeli modelde $R^2=r^2$'dir.
\item[B11.5 A] Yeni bireyin doğal saçılması tahmin aralığına eklenir.
\item[B11.6 A] Veri aralığı dışındaki tahmin ekstrapolasyondur.
\end{description}

\section*{Bölüm 12}
\begin{description}
\item[B12.1 A] Beklenen değer satır ve sütun toplamlarından hesaplanır.
\item[B12.2 A] Serbestlik derecesi $(r-1)(c-1)$'dir.
\item[B12.3 A] Seyrek küçük tablolarda Fisher kesin testi düşünülebilir.
\item[B12.4 A] Cramér $V$ kategorik ilişkinin büyüklüğünü özetler.
\item[B12.5 A] Uygun yüzdeler farklı grup hacimlerini karşılaştırılabilir kılar.
\item[B12.6 A] Standartlaştırılmış artıklar hücre katkılarını gösterir.
\end{description}

\section*{Bölüm 13}
\begin{description}
\item[B13.1 A] Mann--Whitney iki bağımsız sıralı gruba uygundur.
\item[B13.2 A] Wilcoxon eşleştirilmiş farkların sıra büyüklüklerini kullanır.
\item[B13.3 A] Kruskal--Wallis üç veya daha fazla bağımsız grubu karşılaştırır.
\item[B13.4 A] Mann--Whitney her koşulda yalnız medyan testi değildir.
\item[B13.5 A] İşaret testi farkların yalnız pozitif veya negatif yönünü kullanır.
\item[B13.6 A] Sıra testlerinin de bağımsızlık ve tasarım varsayımları vardır.
\end{description}

\section*{Bölüm 14}
\begin{description}
\item[B14.1 A] Yöntem seçimi araştırma sorusundan başlamalıdır.
\item[B14.2 A] Belgelenmeyen elle değişiklikler analizin yeniden üretimini bozar.
\item[B14.3 A] Aynı satır ön--son eşleşmesini korur.
\item[B14.4 A] Olağandışı güçlü sonuç önce veri ve grafiklerle denetlenir.
\item[B14.5 A] Syntax işlem adımlarını yeniden çalıştırılabilir kılar.
\item[B14.6 A] Birincil analiz sonuç görülmeden önce tanımlanmalıdır.
\end{description}